% ==============================================================================
% SECCIÓN 8: ACTIVIDADES Y CRONOGRAMA
% ==============================================================================

\newpage
\section{Actividades y Cronograma}

Las actividades del proyecto se vinculan a los objetivos específicos definidos en el apartado 2 y se desarrollan mediante una metodología ágil. Tras una fase inicial de análisis y especificación de requisitos, que concluye con la entrega del 22/10/2026, el diseño, la implementación y las pruebas se realizan en tres iteraciones, cada una con su entrega, revisión y retrospectiva (RI-03, RI-04, RI-06). Las fechas de cada iteración coinciden con las entregas previstas en el apartado 10, y la planificación detallada de tareas y su seguimiento se registran en JIRA. Los responsables se asignan según los roles del equipo, y la participación del centro se canaliza a través de los cauces descritos en el apartado 8. Cualquier cambio en las fechas o en el reparto de requisitos entre iteraciones se registrará y comunicará al centro.

\subsection{Actividades y entregables previstos}
\begin{center}
\small
\begin{tabularx}{\textwidth}{|>{\raggedright\arraybackslash}p{2.5cm}|>{\raggedright\arraybackslash}p{1.5cm}|>{\raggedright\arraybackslash}X|>{\raggedright\arraybackslash}p{3.4cm}|}
\hline
\rowcolor{tableHeader}
Actividad & Objetivo & Participantes e implicación & Entregable / evidencia \\
\hline
Fase inicial: análisis y especificación & OBJ-01 & Equipo: Coordinador (lidera) y Catalogador (requisitos y backlog). Centro: exposición inicial (25/09), visita al centro (02/10) y dudas por correo. & Requisitos y alcance aprobados (DC-01) \\
\hline
Diseño de flujos y prototipo & OBJ-02 & Equipo: Gestor de Accesibilidad (lidera) y [quien haga Figma]. Centro y usuarios: revisión del prototipo. & Prototipo y observaciones registradas (DC-02) \\
\hline
Implementación del MVP & OBJ-03 & Equipo: [desarrolladores backend y Flutter]. Centro: resolución de dudas. & Versión con requisitos aprobados al final de cada iteración \\
\hline
Pruebas y ajustes & OBJ-04 & Equipo: Gestor de Calidad (lidera). Centro y usuarios: [validación que se acuerde]. & Informe de pruebas e incidencias (DC-03, DC-05) \\
\hline
Documentación y entrega & OBJ-05 & Equipo: Catalogador (manual) y Presentador (presentación final). Centro: revisión y recepción. & Aplicación y documentación final (DC-04, DC-06) \\
\hline
\end{tabularx}
\end{center}

\subsection{Cronograma}
La planificación visual del anexo abarca de octubre a diciembre y se organiza en cuatro sprints (siendo el 0 el sprint actual). La gestión de usuarios y autenticación se sitúa al inicio, la agenda y la gestión de tareas se desarrollan entre los primeros sprints, los pasos de las actividades y los ajustes de accesibilidad se concentran en la parte intermedia y las tareas alternativas aparecen en el último sprint.